\chapter{Заключение}
\label{ch:conclusion}

Цель диссертации состояла в том, чтобы в зафиксированной экспериментальной
области сопоставить исследованные режимы predictive coding с BP, перейти от
конечных метрик к внутреннему механизму вычисления и проверить, может ли
доступное до завершения вычисления состояние использоваться для безопасного и
экономически полезного раннего решения. Для этого работа была организована как
последовательная программа Stage~1/2, Stage~3A, Stage~3B и QWake C1/C2 с
versioned protocols, отдельными admission gates и сохраняемой provenance.

\section{Ответ на RQ1}

RQ1 спрашивал, при каких зафиксированных условиях Exact, FixedPred и Strict
близки к BP и где различия выходят за выбранные границы.

Stage~1/2 показывает, что близость конечного качества может сосуществовать с
различиями вычислительной стоимости. Для зарегистрированного FashionMNIST
paired analysis FixedPred и Strict находятся внутри заранее заданной
equivalence margin по test macro-F1 относительно BP. Stage~2 сохраняет
зарегистрированную final-quality surface при изменившемся runtime-over-BP.

Stage~3A ограничивает интерпретацию такой поведенческой близости. В
зарегистрированной области FixedPred ближе к BP по направлению градиента и
representation metrics, чем Strict, но при этом ранний gradient norm
существенно уменьшен. Поэтому ответ на RQ1 не является бинарным:
``эквивалентен / не эквивалентен''. Близость зависит от выбранного уровня ---
final metric, gradient direction, gradient scale или representation state ---
и должна указываться вместе с соответствующим scope.

В claims contract C01 теперь фиксирует поддержанный Stage~1 paired
macro-F1 equivalence result вместе с сохранением Stage~2 final-quality surface
и cross-version runtime shift; C02 отдельно фиксирует поддержанный Stage~3A
layer-wise результат.

\section{Ответ на RQ2}

RQ2 был посвящён вычислительным механизмам и компонентам стоимости.

Stage~3B B0 локализует \texttt{state\_inference} как основную device-time
область в исследованной ROCm/float32 поверхности; Strict в этой матрице имеет
более высокие device time и peak memory относительно FixedPred. Это
поддерживает C03.

Первичная mechanism-attribution модель SI-MA0 не закрывает COST-MA0 gate, и
этот результат сохраняется как C04. Отдельная SI-MA1 calibration закрывает
CAL-COST-MA1; signed residual интерпретируется как over-closure accounting
model, а не как отрицательная физическая стоимость, что поддерживает C05.

Точные implementation candidates B1 и B2 проходят собственные equivalence
gates до сравнительного профилирования, поддерживая C06. После admission
полностью завершается зарегистрированная 288-cell matched matrix без failures
или retries, что поддерживает C07. При этом диссертация не выводит из
завершённости profiling surface универсальное превосходство какого-либо
candidate.

Ответ на RQ2 поэтому состоит в staged localization: стоимость связывается с
конкретной областью \texttt{state\_inference}, исходная attribution hypothesis
получает отрицательный gate, observer calibration проверяется отдельно, а
альтернативные exact organizations допускаются к matched profiling только
после functional-equivalence checks.

\section{Ответ на RQ3}

RQ3 объединял три последовательные ступени раннего решения: наличие
достаточного для задачи preterminal outcome, возможность безопасно распознать
такой outcome по pre-action состоянию и положительную вычислительную выгоду.

На sealed QWake calibration surface первые две ступени выполняются в bounded
смысле. Наличие безопасно принятых записей означает ненулевую область
post-action oracle-sufficient outcomes, а среди 2\,625 frozen scalar policies
существуют candidates с zero dangerous accepts и ненулевым coverage. Лучшая безопасная policy принимает
216 из 756 records, около 28.57\%, без dangerous accepts. C08 поэтому имеет
статус \texttt{supported}.

Второе условие при frozen full decision-cost accounting не выполняется. Ни
одна из 2\,625 policies не имеет одновременно zero dangerous accepts,
nontrivial coverage и positive aggregate net saving. C09 имеет статус
\texttt{rejected}. Для лучшей безопасной policy полный decision cost
существенно превышает gross safely avoidable suffix, а примерно 99.887\%
accounting cost приходится на observer category.

Этот отрицательный экономический результат не переносится на marginal runtime
cost простого recognizer. Такой estimand не входил в C2, поэтому C10 остаётся
\texttt{not\_tested}. Поскольку eligible policy отсутствовала, C2 не установил
policy freeze и protocol не открыл C3; C11 также остаётся
\texttt{not\_tested}.

Таким образом, RQ3 получает трёхчастный ответ: task-relative sufficient
preterminal outcomes присутствуют на sealed surface; их ненулевое подмножество
безопасно распознаётся внутри frozen family по pre-action признакам; но
safe-and-beneficial policy при frozen full decision-cost accounting не
наблюдается. Marginal implementation viability текущей работой не проверена.

\section{Научный и методологический итог}
\label{sec:conclusion-contribution}

В пределах заявленного scope работа связывает несколько уровней анализа,
которые часто рассматриваются отдельно: конечное поведение, внутренние
градиенты и представления, mechanism-level profiling и bounded decision about
remaining computation. Главный содержательный итог состоит не в выборе одного
``лучшего'' режима, а в локализации того, какие свойства совпадают, какие
различаются и какая следующая гипотеза становится допустимой после каждого
этапа.

Методологически работа использует явное разделение difference и equivalence
claims, model-level statistical units, preservation of negative gates,
pre-action/post-action separation для QWake и admission-by-receipt между
последовательными стадиями. Thesis-facing summaries дополнительно связывают
ключевые численные утверждения с tracked evidence через SHA-256 и проверяются
CI до генерации LaTeX.

Такой подход важен для интерпретации отрицательного результата. SI-MA0 и C2 не
исчезают после появления более узких объяснений. Они остаются ограничениями
соответствующих hypotheses и определяют, какой следующий эксперимент является
научно новым, а какой был бы лишь post-hoc изменением уже завершённого
estimand.

\section{Границы выводов}

Результаты относятся к исследованным версиям Torch2PC, datasets,
\texttt{lenet\_classic}, model seeds, algorithmic regimes, numerical
tolerances и закреплённым hardware/software environments. Они не устанавливают
универсальную эквивалентность PC и BP и не определяют универсально
предпочтительный implementation.

QWake C2 дополнительно ограничен finite candidate family, 756 sealed
calibration records и frozen full decision-cost accounting. Zero dangerous
accepts на этой поверхности не является population-level safety guarantee.
Отсутствие safe-and-beneficial policy не исключает другой recognizer, другую
representation surface или другую marginal implementation architecture.

Эти ограничения не являются внешними оговорками к итоговому результату:
статусы \texttt{supported}, \texttt{rejected} и \texttt{not\_tested}
закреплены в claims contract и используются как границы формулировки выводов.

\section{Последующая работа}

Непосредственный следующий исследовательский вопрос уже локализован
результатами C08--C10: можно ли реализовать найденный или функционально
сопоставимый safe recognizer с настолько низкой \emph{marginal} стоимостью,
чтобы его реальный incremental cost пересёк break-even относительно
предотвращённого suffix computation. Это требует отдельной preregistration,
новой cost surface и независимого измерения; старый C2 не должен
ретроспективно пересчитываться под новую cost model.

После появления admissible marginal-cost candidate возможна независимая
confirmatory campaign на новой evidence surface. Она должна иметь собственную
protocol identity и не рассматривается как задним числом открытый C3 исходной
C2 chain.

Дополнительные направления включают перенос staged analysis на другие
архитектуры, datasets, precision regimes и hardware backends, а также
исследование более дешёвых observer representations. Во всех случаях
существенно сохранять разделение поведенческого сходства, внутреннего
механизма, safety и экономической стоимости.

\section{Итоговый вывод}

Завершённая экспериментальная программа последовательно сузила исходный
вопрос о связи predictive coding и backpropagation. В исследованной области
близость конечного качества не означает идентичности градиентов,
представлений или runtime; mechanism-level profiling локализует существенную
стоимость внутри \texttt{state\_inference}; а QWake показывает, что наличие
безопасно распознаваемого pre-action сигнала само по себе не устанавливает
экономическую целесообразность его использования.

Именно разделение этих уровней составляет итог работы: informational
feasibility может быть поддержана при одновременно отклонённой
frozen-accounting feasibility, а marginal implementation feasibility при этом
остаётся открытым, но уже точно сформулированным следующим вопросом. На этом
завершённая evidence chain достигает своего preregistered stopping boundary:
C2 остаётся bounded-negative по economic gate, C3 не открывается, и новые
утверждения требуют нового экспериментального контракта.
